\section{Related Work}

\textbf{Purpose-based access control.} PBAC binds access permissions to the
purpose of a data request~\cite{byun2008}; variants attach privacy-aware
purpose constraints at storage and query time. Usage control extends this with
obligations and conditions at data-consumption time~\cite{usagecontrol2018}.
These systems govern \emph{release}: they decide who may read what, for which
declared reason. They do not observe the downstream decision, and they cannot
detect a model weighting a field it was permitted to read for one purpose while
acting under another.

\textbf{Contextual integrity.} Nissenbaum's contextual integrity treats
appropriateness of data flow as the central privacy criterion~\cite{nissenbaum2010},
and Barth et al.\ derive a formal language for context-relative norms~\cite{barth2006}.
CI-Work operationalizes flow norms as runtime checks for
applications~\cite{ciwork2026}. These lines treat \emph{flow} as the regulated
object. Our benchmark treats \emph{influence} as the regulated object: a flow
can be perfectly legal and still causally tilt a prohibited decision.

\textbf{LLM privacy benchmarks.} ToolPrivacyBench audits whether agent tools
leak or mishandle private data in tool calls~\cite{toolprivacybench2026};
CONFAIDE generates contextualized private-data scenarios~\cite{confaide2024};
PrivaCI-Bench builds role-negotiation tasks for privacy context
inference~\cite{privacibench2025}. These benchmarks ask whether private data
is \emph{exposed} to a downstream consumer or tool. FinBoundBench instead asks
whether a decision \emph{changes} when the data's purpose flips---the
transmission may be identical.

\textbf{Fairness and bias influence.} Counterfactual fairness asks whether
predictions are invariant under counterfactual changes of protected
attributes~\cite{kusner2017}. Fin-Bias measures which prompt features drive
financial-advice behavior under instruction changes~\cite{finbias2026}.
Both are methodologically close to us: they vary an input while holding the
case fixed. But they target \emph{bias} (protected groups, advice consistency),
not \emph{purpose authorization}; neither reports against a decision
floor, and neither is preregistered with frozen artifacts. IFC-ML analyzes
information-flow control for ML pipelines~\cite{ifcml2024} at the system level
rather than the decision level.

\textbf{Regulatory framing.} GDPR Articles 5(1)(b) and 5(1)(c) ground purpose
limitation and minimization~\cite{gdpr2016}; the EU AI Act Article 6 Annex III
classifies creditworthiness and insurance-risk systems as high-risk, tying
purpose discipline to deployment obligations~\cite{euaiact2024}. Work on
reviving purpose limitation in data-driven systems~\cite{purposeminimization2021}
argues that technical enforcement lags legal obligation---the gap this paper
addresses with a measurable property.

\textbf{Position.} The closest conceptual relative is ToolPrivacyBench
(transmission of private data) and the closest methodological relative is
Fin-Bias (input perturbation). To our knowledge no existing benchmark measures,
on the same confidential signal, the \emph{joint} requirement of authorized
utility retention and floor-relative prohibited influence, under a
preregistered protocol with frozen, independently re-computed evidence.
